% !TEX root = ../main.tex
\chapter{Spécification des besoins et analyse}

% ─────────────────────────────────────────────────────────────
\section{Introduction}

Après avoir posé le cadre général du projet et défini la méthodologie de travail adoptée,
nous entamons dans ce chapitre la phase d'analyse fonctionnelle d'ASM Track. Cette étape
consiste à identifier les acteurs du système, à recenser leurs besoins fonctionnels et non
fonctionnels, puis à les formaliser sous forme de diagrammes de cas d'utilisation et de
diagrammes de séquence. Cette analyse constitue le fondement sur lequel reposent tous
les choix de conception et d'implémentation qui seront détaillés dans les chapitres suivants.

% ─────────────────────────────────────────────────────────────
\section{Étude fonctionnelle}

L'étude fonctionnelle est une phase indispensable pour identifier les tâches que le système
doit accomplir et les interactions entre ses différents utilisateurs. Elle se traduit par la
liste des acteurs, leurs besoins et leurs interactions, formalisés sous forme de diagrammes UML.

\subsection{Identification des acteurs}

Un acteur représente toute entité externe qui interagit directement avec le système.
Nous avons identifié quatre acteurs pour ASM Track :

\begin{itemize}
  \item \textbf{Le Super Admin} : représente l'équipe d'ASM. Il dispose d'un accès
        global à la plateforme, sans restriction par entreprise. Il est responsable
        de la gestion des entreprises clientes, des chauffeurs et des utilisateurs
        administrateurs.

  \item \textbf{L'Admin / Dispatcher} : représente un utilisateur appartenant à une
        entreprise cliente. Il gère les livraisons, planifie les tournées, suit
        l'avancement en temps réel et génère les rapports. Ses données sont
        strictement isolées de celles des autres entreprises.

  \item \textbf{Le Chauffeur} : représente un livreur sur le terrain. Il interagit
        exclusivement via l'application mobile Flutter. Il consulte ses livraisons
        assignées, confirme les étapes du cycle de livraison et capture les preuves
        de livraison, même sans connexion internet.

  \item \textbf{Le Système ERP} : représente le système de gestion de l'entreprise
        cliente (Odoo ou Dux). C'est un acteur externe qui envoie des commandes
        vers ASM Track et reçoit en retour les mises à jour de statut des livraisons.
\end{itemize}

\subsection{Identification des besoins fonctionnels}

\subsubsection*{En tant que Super Admin, je peux :}

\begin{itemize}
  \item Créer, modifier, activer et désactiver des entreprises clientes.
  \item Créer et gérer les profils des chauffeurs (nom, téléphone, statut).
  \item Créer et gérer les comptes administrateurs des entreprises clientes.
  \item Réinitialiser les mots de passe des administrateurs.
  \item Consulter les statistiques globales de la plateforme.
  \item Accéder aux données de toutes les entreprises sans restriction de tenant.
\end{itemize}

\subsubsection*{En tant qu'Admin / Dispatcher, je peux :}

\begin{itemize}
  \item M'authentifier avec mon email et mon mot de passe.
  \item Consulter et rechercher les commandes non planifiées par statut, date ou source.
  \item Replanifier une commande en la retirant d'une tournée et en la réassignant
        à une autre.
  \item Annuler une commande avant qu'elle ne soit prise en charge par un chauffeur.
  \item Accéder au constructeur de tournées et visualiser les commandes non planifiées
        disponibles.
  \item Créer une tournée et y ajouter des commandes par glisser-déposer.
  \item Optimiser l'itinéraire de la tournée automatiquement via OSRM.
  \item Assigner une tournée à un chauffeur disponible.
  \item Suivre en temps réel l'avancement des tournées et la position des chauffeurs.
  \item Recevoir des alertes automatiques en cas de dépassement de délai SLA.
  \item Consulter les preuves de livraison (photos et signatures) capturées par
        les chauffeurs lors de l'exécution de la tournée.
  \item Générer un rapport PDF d'une tournée clôturée avec ses indicateurs de performance.
  \item Configurer les seuils SLA de l'entreprise (délai d'attente, délai d'assignation,
        délai de ramassage) et les modifier dynamiquement en temps réel.
\end{itemize}

\subsubsection*{En tant que Chauffeur, je peux :}

\begin{itemize}
  \item M'authentifier avec mon numéro de téléphone et mon mot de passe.
  \item Consulter ma tournée du jour avec la liste ordonnée des arrêts à effectuer.
  \item Démarrer ma tournée et naviguer vers chaque arrêt grâce à l'itinéraire calculé.
  \item Confirmer le ramassage des colis d'un arrêt (statut \texttt{PICKED\_UP}).
  \item Démarrer le transit vers le destinataire (statut \texttt{IN\_TRANSIT}).
  \item Confirmer la livraison complète d'un arrêt avec capture de photo et de signature (POD).
  \item Confirmer une livraison partielle lorsqu'une partie seulement des articles
        a été remise au destinataire (statut \texttt{PARTIALLY\_DELIVERED}).
  \item Marquer un arrêt comme échoué avec indication de motif.
  \item Travailler en mode hors-ligne : toutes mes actions sur la tournée sont
        sauvegardées localement et synchronisées automatiquement dès le retour du réseau.
  \item Transférer ma tournée à un autre chauffeur via QR code.
  \item Mettre à jour ma position GPS pour permettre le suivi en temps réel par le dispatcher.
\end{itemize}

\subsubsection*{En tant que Système ERP, je peux :}

\begin{itemize}
  \item Envoyer des commandes confirmées vers ASM Track pour les transformer en livraisons.
  \item Recevoir automatiquement les mises à jour de statut des livraisons (livrée,
        partiellement livrée, échouée, annulée) ainsi que les ajustements de stock
        correspondants, sans intervention manuelle.
\end{itemize}

\subsection{Identification des besoins non fonctionnels}

Les besoins non fonctionnels définissent les critères de qualité que le système doit
respecter indépendamment des fonctionnalités :

\begin{itemize}
  \item \textbf{La sécurité} : toutes les communications sont protégées par HTTPS.
        L'authentification repose sur un serveur OAuth2 personnalisé émettant des tokens
        JWT signés avec une clé RSA. L'API Gateway valide chaque token avant de router
        la requête. L'isolation des données entre entreprises est assurée par des filtres
        Hibernate appliqués automatiquement.

  \item \textbf{La performance} : les mises à jour temps réel doivent être poussées en
        moins de 2 secondes. Les pages du tableau de bord doivent se charger en moins
        de 3 secondes dans des conditions réseau normales.

  \item \textbf{La disponibilité} : la plateforme doit être disponible en continu.
        Chaque microservice dispose d'un healthcheck Docker et redémarre automatiquement
        en cas de défaillance.

  \item \textbf{La fiabilité} : aucun événement métier ne doit être perdu. La
        synchronisation ERP est garantie même en cas de panne réseau ou de redémarrage
        d'un service.

  \item \textbf{L'idempotence} : les opérations critiques telles que la synchronisation
        ERP et la soumission de preuves de livraison (POD) peuvent être rejouées sans
        effet de bord. Une même requête exécutée plusieurs fois produit toujours le
        même résultat, évitant ainsi les doublons en cas de retry automatique.

  \item \textbf{La scalabilité} : l'architecture microservices permet de dimensionner
        indépendamment chaque service selon la charge. Le modèle multi-tenant permet
        d'intégrer de nouveaux clients sans déploiement supplémentaire.

  \item \textbf{La maintenabilité} : le code métier est isolé des dépendances externes
        grâce à une architecture en couches où chaque système tiers (ERP, stockage, routage)
        est remplaçable sans modifier le reste de la plateforme.

  \item \textbf{L'accessibilité mobile} : l'application chauffeur doit fonctionner
        avec ou sans connexion internet. Les données critiques sont stockées localement
        et synchronisées automatiquement au retour du réseau.
\end{itemize}

\subsection{Diagrammes de cas d'utilisation}

Le diagramme de cas d'utilisation décrit le comportement du système du point de vue
de l'utilisateur. Il exprime les fonctionnalités offertes par le système et les interactions
entre les acteurs et ces fonctionnalités.

\subsubsection*{Diagramme de cas d'utilisation global}

La figure \ref{fig:uc-global} présente le diagramme de cas d'utilisation global d'ASM Track.
Il illustre l'ensemble des interactions entre les quatre acteurs du système — Super Admin,
Admin/Dispatcher, Chauffeur et Système ERP — et les fonctionnalités qu'ils peuvent réaliser.
Les relations \textit{<<include>>} indiquent les cas d'utilisation obligatoirement exécutés
(notamment l'authentification), tandis que les relations \textit{<<extend>>} représentent
les fonctionnalités optionnelles déclenchées sous certaines conditions.

\begin{figure}[H]
  \centering
  \includegraphics[width=\textwidth]{globalDigram}
  \caption{Diagramme de cas d'utilisation global — ASM Track}
  \label{fig:uc-global}
\end{figure}

\subsubsection*{Diagramme raffiné — Création de tournée}

La figure suivante présente le diagramme raffiné de la création de tournée,
illustrant les actions disponibles pour le dispatcher lors de la planification.

\begin{figure}[H]
  \centering
  \includegraphics[width=\textwidth]{CreateRoute}
  \caption{Diagramme raffiné — Création de tournée}
  \label{fig:uc-creation-tournee}
\end{figure}

\subsubsection*{Diagramme raffiné — Gestion et suivi des tournées}

La figure suivante présente le diagramme raffiné de la gestion et du suivi
des tournées, couvrant les actions du dispatcher lors de l'exécution en temps réel.

\begin{figure}[H]
  \centering
  \includegraphics[width=\textwidth]{GererRoute}
  \caption{Diagramme raffiné — Gestion et suivi des tournées}
  \label{fig:uc-gestion}
\end{figure}

\subsubsection*{Diagramme raffiné — Exécution et confirmation de livraison}

La figure suivante présente le diagramme raffiné du côté chauffeur, illustrant
le cycle d'exécution d'une tournée et les actions de confirmation de livraison.

\begin{figure}[H]
  \centering
  \includegraphics[width=\textwidth]{UcDriver}
  \caption{Diagramme raffiné — Exécution et confirmation de livraison}
  \label{fig:uc-driver}
\end{figure}

\textbf{Description textuelle CU : Assigner une livraison}

\begin{table}[H]
\centering
\begin{tabularx}{\textwidth}{|L{3.5cm}|X|}
\hline
\rowcolor{primary!10}
\multicolumn{2}{|c|}{\textbf{CU : Assigner une livraison à un chauffeur}} \\
\hline
\textbf{Acteur} & Admin / Dispatcher \\
\hline
\textbf{Précondition} & L'admin est authentifié ; la livraison est en statut \texttt{UNSCHEDULED} \\
\hline
\multicolumn{2}{|l|}{\textit{Enchaînement principal :}} \\
\hline
\multicolumn{2}{|X|}{
  \begin{enumerate}[leftmargin=1.5cm]
    \item L'admin sélectionne une livraison dans le tableau de bord.
    \item Le système affiche les chauffeurs disponibles.
    \item L'admin choisit un chauffeur et confirme l'assignation.
    \item Le système met à jour le statut à \texttt{SCHEDULED} et envoie
          une notification push FCM au chauffeur.
    \item Le système enregistre l'événement dans l'historique des statuts.
  \end{enumerate}
} \\
\hline
\textbf{Postcondition} & La livraison est assignée ; le chauffeur reçoit une notification sur son mobile \\
\hline
\textbf{Enchaînement alternatif} &
  Si aucun chauffeur n'est disponible, le système affiche un message d'avertissement
  et l'admin peut choisir de planifier l'assignation ultérieurement. \\
\hline
\end{tabularx}
\caption{Description textuelle CU : Assigner une livraison}
\end{table}

\subsubsection*{Diagramme raffiné — Planification des tournées}

\begin{figure}[H]
  \centering
  \includegraphics[width=\textwidth]{CreateRoute}
  \caption{Diagramme raffiné — Création de tournée}
  \label{fig:uc-tournees}
\end{figure}

\textbf{Description textuelle CU : Créer une tournée}

\begin{longtable}{|p{3.5cm}|p{11.5cm}|}
\caption{Description textuelle CU : Créer une tournée} \\
\hline
\rowcolor{primary!10}
\multicolumn{2}{|c|}{\textbf{CU : Créer une tournée}} \\
\hline
\textbf{Acteur} & Admin / Dispatcher \\
\hline
\textbf{Précondition} &
  L'admin est authentifié sur le tableau de bord ; des commandes en statut
  \texttt{UNSCHEDULED} existent ; au moins un chauffeur, un véhicule et un
  dépôt sont disponibles. \\
\hline
\multicolumn{2}{|l|}{\textit{Enchaînement principal :}} \\
\hline
\multicolumn{2}{|p{15.4cm}|}{
  \begin{enumerate}[leftmargin=1.5cm]
    \item L'admin ouvre le constructeur de tournées et clique sur
          \textit{"Nouvelle tournée"}.
    \item Le système affiche un formulaire de création. L'admin saisit le nom
          de la tournée, la date d'opération, sélectionne le chauffeur, le
          véhicule et le dépôt de départ.
    \item Le système crée la tournée en statut \texttt{DRAFT} et affiche
          le tableau des commandes non planifiées.
    \item L'admin glisse-dépose les commandes depuis le tableau ou depuis les
          épingles de la carte vers la tournée.
    \item Le système ajoute chaque commande comme arrêt et met à jour
          le compteur de charge (poids total / capacité véhicule).
    \item L'admin clique sur \textit{"Optimiser"} pour calculer l'itinéraire
          optimal via le moteur OSRM.
    \item Le système retourne une suggestion d'ordre des arrêts avec les
          distances, durées, fenêtres horaires optimisées et l'heure
          d'arrivée estimée (ETA) pour chaque arrêt.
    \item L'admin consulte la suggestion et clique sur \textit{"Appliquer"}.
    \item Le système enregistre l'ordre optimisé et les fenêtres horaires
          suggérées pour chaque arrêt.
  \end{enumerate}
} \\
\hline
\textbf{Postcondition} &
  La tournée est créée en statut \texttt{DRAFT} avec les arrêts ordonnés,
  les fenêtres horaires calculées et les ETAs enregistrés. \\
\hline
\textbf{Enchaînement alternatif} &
  \textbf{A1 — Véhicule en surcharge :} Si le poids total des commandes
  dépasse la capacité du véhicule, le système affiche une alerte rouge.
  L'admin peut retirer des commandes ou changer de véhicule.

  \textbf{A2 — OSRM indisponible :} Si le moteur de routage ne répond pas,
  le système conserve l'ordre manuel des arrêts et affiche un avertissement.
  L'admin peut relancer l'optimisation ultérieurement. \\
\hline
\end{longtable}

\textbf{Description textuelle CU : Valider une tournée}

\begin{longtable}{|p{3.5cm}|p{11.5cm}|}
\caption{Description textuelle CU : Valider une tournée} \\
\hline
\rowcolor{primary!10}
\multicolumn{2}{|c|}{\textbf{CU : Valider une tournée}} \\
\hline
\textbf{Acteur} & Admin / Dispatcher \\
\hline
\textbf{Précondition} &
  L'admin est authentifié ; la tournée est en statut \texttt{DRAFT} ;
  tous les arrêts disposent de fenêtres horaires complètes ; aucune
  violation chronologique n'est détectée. \\
\hline
\multicolumn{2}{|l|}{\textit{Enchaînement principal :}} \\
\hline
\multicolumn{2}{|p{15.4cm}|}{
  \begin{enumerate}[leftmargin=1.5cm]
    \item L'admin ouvre la tournée dans le constructeur et clique sur
          \textit{"Finaliser et valider"}.
    \item Le système ouvre une fenêtre de confirmation affichant pour chaque
          arrêt : le numéro, le client, l'adresse, les horaires prévus
          et un indicateur vert ou rouge selon la validité des fenêtres.
    \item Le système vérifie automatiquement que chaque arrêt commence
          après la fin du précédent et qu'aucun chevauchement n'existe.
    \item Si toutes les vérifications sont satisfaites, le bouton
          \textit{"Valider la tournée"} est activé et affiché en vert.
    \item L'admin clique sur \textit{"Valider la tournée"}.
    \item Le système sauvegarde les fenêtres horaires finales et fait
          passer le statut de la tournée de \texttt{DRAFT} à
          \texttt{VALIDATED}.
    \item Le système envoie une notification push FCM au chauffeur
          assigné pour l'informer que sa tournée est prête.
  \end{enumerate}
} \\
\hline
\textbf{Postcondition} &
  La tournée est en statut \texttt{VALIDATED}, les fenêtres horaires
  sont verrouillées et le chauffeur a reçu une notification. \\
\hline
\textbf{Enchaînement alternatif} &
  \textbf{A1 — Violations chronologiques détectées :} Si un ou plusieurs
  arrêts présentent des chevauchements ou des incohérences d'horaires,
  le système les surligne en rouge et désactive le bouton de validation.
  L'admin doit corriger les fenêtres horaires concernées avant de
  pouvoir valider.

  \textbf{A2 — Fenêtres horaires incomplètes :} Si un arrêt ne possède
  pas de fenêtre horaire définie, le système bloque la validation et
  indique les arrêts manquants. \\
\hline
\end{longtable}

\subsubsection*{Diagramme raffiné — Exécution et confirmation de livraison}

\begin{figure}[H]
  \centering
  \includegraphics[width=\textwidth]{UcDriver}
  \caption{Diagramme raffiné — Exécution et confirmation de livraison (Chauffeur)}
  \label{fig:uc-driver-2}
\end{figure}

\textbf{Description textuelle CU : Confirmer une livraison avec capture de preuve (POD)}

\begin{longtable}{|p{3.5cm}|p{11.5cm}|}
\caption{Description textuelle CU : Confirmer une livraison avec POD} \\
\hline
\rowcolor{primary!10}
\multicolumn{2}{|c|}{\textbf{CU : Confirmer une livraison avec capture de preuve (POD)}} \\
\hline
\textbf{Acteur} & Chauffeur \\
\hline
\textbf{Précondition} &
  Le chauffeur est authentifié sur l'application mobile ; la tournée
  est en statut \texttt{IN\_PROGRESS} ; l'arrêt concerné est en
  statut \texttt{IN\_TRANSIT}. \\
\hline
\multicolumn{2}{|l|}{\textit{Enchaînement principal :}} \\
\hline
\multicolumn{2}{|p{15.4cm}|}{
  \begin{enumerate}[leftmargin=1.5cm]
    \item Le chauffeur consulte sa tournée du jour et sélectionne
          l'arrêt en cours dans l'application mobile.
    \item Le chauffeur démarre le transit vers l'arrêt. Le système
          met à jour le statut de l'arrêt à \texttt{IN\_TRANSIT}.
    \item À l'arrivée chez le destinataire, le chauffeur ouvre le
          formulaire de confirmation de livraison.
    \item Le chauffeur imprime le bon de livraison (BL) et le fait
          signer par le destinataire.
    \item Le chauffeur capture deux photos via l'application mobile :
          une photo du colis remis et une photo du bon de livraison signé.
    \item L'application envoie les deux photos au serveur comme
          preuve de livraison (POD).
    \item Le serveur stocke les photos sur MinIO et génère des URLs
          d'accès sécurisées.
    \item Le système met à jour le statut de l'arrêt à \texttt{DELIVERED}
          et enregistre l'horodatage de livraison.
    \item Le système publie un événement sur RabbitMQ. L'adaptateur ERP
          correspondant synchronise automatiquement la livraison dans
          le système de l'entreprise (Odoo ou Dux).
    \item Le dispatcher voit instantanément le statut mis à jour
          sur son tableau de bord, sans avoir à recharger la page.
  \end{enumerate}
} \\
\hline
\textbf{Postcondition} &
  L'arrêt est en statut \texttt{DELIVERED}, la preuve de livraison
  est stockée sur MinIO, l'ERP est synchronisé et le dispatcher
  est notifié en temps réel. \\
\hline
\textbf{Enchaînement alternatif} &
  \textbf{A1 — Mode hors-ligne :} Si le chauffeur n'a pas de connexion
  internet, l'application stocke la POD localement dans la file
  d'attente Hive. Dès que la connexion est rétablie, la synchronisation
  s'effectue automatiquement (3 tentatives maximum, TTL 24h).

  \textbf{A2 — Livraison partielle :} Si une partie seulement des
  articles a été remise, le chauffeur sélectionne les articles livrés
  et confirme une livraison partielle. Le statut passe à
  \texttt{PARTIALLY\_DELIVERED} et l'ERP est notifié des quantités
  réellement livrées. \\
\hline
\end{longtable}

\subsection{Diagrammes de séquence système}

Les diagrammes de séquence modélisent les interactions entre les composants du système
lors de l'exécution d'un cas d'utilisation précis. Ils illustrent l'ordre chronologique des
échanges de messages entre les acteurs et les services.

\subsubsection*{Diagramme de séquence — Authentification admin}

\begin{figure}[H]
  \centering
  \includegraphics[width=\textwidth]{Sequencelogin}
  \caption{Diagramme de séquence — Authentification administrateur}
  \label{fig:seq-auth-admin}
\end{figure}

Lorsque l'administrateur soumet ses identifiants, le \textit{UserManagementMS} vérifie
d'abord l'existence et l'état actif du compte en base de données, puis délègue la
validation du mot de passe au serveur d'autorisation via un \textit{password grant} OAuth2.
Au démarrage du service, la clé publique RSA est récupérée depuis l'endpoint JWKS et
mise en cache mémoire, ce qui permet de vérifier la signature des tokens sans aller-retour
réseau. En cas de succès, le token JWT (RS256, 15 min) et le profil de l'utilisateur sont
retournés et stockés dans des cookies \texttt{HttpOnly}.

\subsubsection*{Diagramme de séquence — Synchronisation ERP}

\begin{figure}[H]
  \centering
  \includegraphics[width=\textwidth]{syncSequence}
  \caption{Diagramme de séquence — Synchronisation ERP via patron Outbox}
  \label{fig:seq-outbox-erp}
\end{figure}

Lorsque le chauffeur confirme une livraison, le \textit{DeliveryMS} écrit la mise à jour
du statut et l'événement de synchronisation dans la \textbf{même transaction PostgreSQL},
garantissant l'atomicité. Un \texttt{OutboxProcessor} interroge la table toutes les 20
secondes via \texttt{FOR UPDATE SKIP LOCKED} — technique qui permet à plusieurs instances
de traiter des lots différents sans collision. L'\textit{ErpAdapterMS} reçoit l'événement,
vérifie l'idempotence via une table H2 en mémoire, puis appelle le système ERP cible.
En cas d'indisponibilité, le compteur \texttt{retry\_count} est incrémenté et l'événement
repasse en \texttt{PENDING} jusqu'à 50 tentatives, après quoi une alerte webhook est déclenchée.

% ─────────────────────────────────────────────────────────────
\section{Conclusion}

Ce chapitre a permis de définir précisément le périmètre fonctionnel d'ASM Track.
Nous avons identifié quatre acteurs et décrit leurs besoins fonctionnels sous forme
de listes et de diagrammes de cas d'utilisation avec leurs descriptions textuelles.
Les besoins non fonctionnels fixent les exigences de qualité en matière de sécurité,
de performance et de fiabilité. Les diagrammes de séquence illustrent les flux
d'interaction clés du système. Le chapitre suivant présente l'architecture technique
retenue pour répondre à l'ensemble de ces besoins.
